% from ega4-2-fr.tex
\oldpage[IV-2]{175}
II) \emph{Réduction au cas où $B$ est une $A$-algèbre intègre de type fini
contenant $A$.} --- Soit $(B_\alpha)$ la famille filtrante croissante des
sous-$A$-algèbres de type fini de $B$, et soit $C_\alpha$ la fermeture
intégrale de $A$ dans $B_\alpha$. Il résulte aussitôt de la définition de
la fermeture intégrale que $C$ est réunion des $C_\alpha$; si l'on pose
$B'_\alpha=B_\alpha\otimes_A A'$, $C'_\alpha=C_\alpha\otimes_A A'$,
$C'_\alpha$ est contenu dans $B'_\alpha$ par platitude, et pour la même
raison $B'$ est réunion de la famille filtrante croissante des
$B'_\alpha$ et $C'$ réunion de la famille filtrante croissante des
$C'_\alpha$. Si l'on prouve que $C'_\alpha$ est la fermeture intégrale de
$A'$ dans $B'_\alpha$, il en résultera aussitôt que $C'$ est la fermeture
intégrale de $A'$ dans $B'$. On peut donc se borner au cas où $B$ est une
$A$-algèbre de type fini, donc noethérienne; soient $\mathfrak q_i$
$(1\leqslant i\leqslant n)$ ses idéaux premiers minimaux; comme $B$ est
supposé réduit, il s'identifie à un sous-anneau du produit $B_0$ des
$B/\mathfrak q_i$; si $C_0$ est la fermeture intégrale de $A$ dans $B_0$,
on a $C=B\cap C_0$; si l'on pose $B'_0=B_0\otimes_A A'$,
$C'_0=C_0\otimes_A A'$, on a, par platitude,
$C'=B'\cap C'_0$ (0$_{\mathrm I}$, 6.1.3); il suffit donc de prouver que
$C'_0$ est la fermeture intégrale de $A'$ dans $B'_0$. Mais $C_0$ est
composé direct des $C_i$, fermetures intégrales de $A$ dans les
$B_i=B/\mathfrak q_i$; par suite $C'_0$ est composé direct des
$C'_i=C_i\otimes_A A'$ et il suffit de montrer que $C'_i$ est la fermeture
intégrale de $A'$ dans $B'_i=B_i\otimes_A A'$. On est ainsi ramené au cas
où $B$ est intègre et une $A$-algèbre de type fini; si $\mathfrak p$ est le
noyau de l'homomorphisme $A\to B$, on a aussi
$B'=B\otimes_{A/\mathfrak p}(A'/\mathfrak pA')$; comme le morphisme
$\operatorname{Spec}(A'/\mathfrak pA')\to
\operatorname{Spec}(A/\mathfrak p)$ est normal, on peut remplacer $A$ par
$A/\mathfrak p$ et $A'$ par $A'/\mathfrak pA'$, et supposer par suite que
$A\subset B$.

III) \emph{Cas où $A$ est un corps, $B$ un corps, extension de type fini de
$A$ et tel que $A$ soit algébriquement fermé dans $B$.} --- On a alors
$C=A$, et $A'$ est une $A$-algèbre noethérienne géométriquement normale,
donc composée directe d'anneaux intègres et intégralement clos $A'_i$, les
corps des fractions $K'_i$ des $A'_i$ étant des extensions séparables de
$A$ (4.6.1). La $A'$-algèbre $B'$ est donc composée directe des
$B\otimes_A A'_i=B'_i$ et l'on peut donc se borner au cas où $A'$ est
intègre, son corps des fractions $K'$ étant une extension séparable de $A$.
Or, comme $B$ est un $A$-module plat, $B\otimes_A A'=B'$ s'identifie à un
sous-anneau de $B\otimes_A K'$; comme $K'$ est extension séparable de $A$,
$B\otimes_A K'$ est réduit (4.3.7), et comme de plus $B$ est une extension
primaire de $A$, $B\otimes_A K'$ est intègre (4.3.2); il en est par suite
de même de $B'$. Soit $L'$ le corps des fractions de $B\otimes_A K'$ (qui
est aussi celui de $B'$); c'est évidemment un corps composé $B(K')$ de ses
sous-corps $B$ et $K'$; comme $A$ est algébriquement fermé dans $B$, que
$K'$ est une extension séparable de $A$, et que $B$ et $K'$ sont
linéairement disjoints sur $A$, $K'$ est algébriquement fermé dans
$B(K')=L'$ (Bourbaki, Alg., chap. V, § 9, exerc. 2); a fortiori $A'$, étant
intégralement clos, est intégralement fermé dans $L'$, donc dans $B'$, ce
qui prouve la proposition dans ce cas.

IV) \emph{Cas où $A$ est intègre, $B$ un corps, extension finie du corps
des fractions $K$ de $A$.} --- Alors
$B'=(A'\otimes_A K)\otimes_K B$ est une $B$-algèbre noethérienne
géométriquement normale, et par suite $B'$ est composé direct d'anneaux
intègres; l'anneau total des fractions $L'$ de $B'$ est alors composé
direct d'un nombre fini de corps. Cela étant, comme $B$ est le corps des
fractions de $C$, $B'$ est un anneau de fractions de $C'$ et $L'$ est donc
aussi l'anneau total des fractions de $C'$. Or (6.14.1) comme $C$ est un
anneau normal et $\operatorname{Spec}(A')\to\operatorname{Spec}(A)$ un
morphisme normal d'anneaux noethériens, $C'$ est un anneau normal; mais il
résulte de (6.14.1.1) que $C'$ est alors intégralement fermé dans $L'$, et
a fortiori dans $B'$, et par suite est la fermeture intégrale de $A'$ dans
$B'$.

